\documentclass[journal]{IEEEtran}

\usepackage{cite}
\usepackage{amsmath,amssymb,amsfonts}
\usepackage{algorithmic}
\usepackage{graphicx}
\usepackage{textcomp}
\usepackage{xcolor}
\usepackage{booktabs}
\usepackage{multirow}
\usepackage{hyperref}

% ---------------------------------------------------------------------------
% Draft-only helper: visually flags unfinished content so nothing pending is
% mistaken for a finished claim. Remove this package/command before final
% submission (do NOT leave \todo markers in a submitted manuscript).
% ---------------------------------------------------------------------------
\newcommand{\todo}[1]{\textcolor{red}{\textbf{[TODO: #1]}}}

\begin{document}

\title{DH$^2$A-KT: Dynamic Heterogeneous Hypergraph and Agentic Knowledge
Tracing with Leakage-Controlled Graph Construction}

\author{Tuan~Dao~Minh, Khanh-Trinh~Nguyen, Duong~Nguyen~Tien, Quoc~Khanh~Ngo,
Van-Hau~Nguyen, and~Le~Hoang~Son
\thanks{T. Dao Minh, K.-T. Nguyen, D. N. Tien, and V.-H. Nguyen are with the
Department of Quality Assurance and Testing, Hung Yen University of
Technology and Education, Hung Yen, Vietnam (e-mail: tuanymc@utehy.edu.vn;
trinhnk@utehy.edu.vn; duongnt@utehy.edu.vn; nvhau66@gmail.com).}
\thanks{Q. K. Ngo is with the Advanced International Research Center on
Applied Artificial Intelligence, Information Technology Institute, Vietnam
National University, Hanoi, Vietnam (e-mail: quockhanhngo.official@gmail.com).}
\thanks{L. H. Son is with the Information Technology Institute, Vietnam
National University, 144 Xuan Thuy Street, Cau Giay, Hanoi, Vietnam. He is
the corresponding author (e-mail: sonlh@vnu.edu.vn).
\todo{Manuscript received/revised dates -- fill at submission time.}}}

\markboth{IEEE Transactions on Learning Technologies}%
{Dao \MakeLowercase{\textit{et al.}}: DH$^2$A-KT}

\maketitle

\begin{abstract}
Graph-augmented knowledge tracing (KT) models student mastery by propagating
information over auxiliary concept graphs, but two gaps limit current
practice. First, most graph-KT architectures represent only pairwise
knowledge-component relations and omit heterogeneous entities -- teachers,
hints, forum discussions, instructional videos -- that plausibly shape a
learner's trajectory. Second, when graph-augmented KT models are evaluated
without auditing whether the auxiliary graph was itself built without
leakage, reported accuracy gains are difficult to attribute to genuine
structural signal versus held-out information reaching the model through the
graph channel. We present DH$^2$A-KT, a two-tier architecture that addresses
both gaps while remaining computationally tractable on a single consumer
GPU. Tier 1 (DH$^2$-KT) extends a leakage-controlled, train-only concept
graph construction and audit protocol -- released with a companion paper,
P0 -- from pairwise prerequisite/similarity edges to heterogeneous
hyperedges (session, concept-prerequisite, discussion-thread,
teacher-intervention), and adds a propensity-score causal-effect layer to
separate genuine intervention effects from spurious correlation. We extend
P0's four-diagnostic leakage taxonomy with three hyperedge-specific leakage
classes -- group-membership, cross-modal, and intervention-timing leakage --
that a pairwise audit cannot detect. Tier 2 (AgentDG-KT) is a
deliberately-scoped pilot in which a small team of LLM agents
(Diagnostician, Tutor/Hint, Critic) reads Tier 1's frozen output to produce
natural-language explanations, gated by a mandatory Critic agent rather than
benchmarked at the same quantitative scale as Tier 1. We reuse P0's already
audited, leakage-controlled baseline results on XES3G5M, ASSISTments 2012,
and Junyi Academy as the Tier 1 comparison column, avoiding redundant
retraining of nine baseline knowledge-tracing models.
\todo{Replace this sentence with the actual Tier 1/Tier 2 headline result
once M5/M6/M9 (see docs/execution-plan.md) are complete -- do not
pre-announce a number before it exists.}
\end{abstract}

\begin{IEEEkeywords}
Knowledge tracing, hypergraph neural networks, heterogeneous graphs, large
language model agents, causal inference, data leakage, educational data
mining.
\end{IEEEkeywords}

% =============================================================================
\section{Introduction}
% =============================================================================

Knowledge tracing (KT) estimates a learner's latent mastery of knowledge
components (KCs) from interaction history, supporting adaptive practice,
intelligent tutoring, and curriculum recommendation. The field has
progressed from probabilistic and sequence models
\cite{piech2015dkt} to graph-augmented models that propagate information
over auxiliary concept graphs \cite{nakagawa2019gkt,yang2021gikt}, and, most
recently, toward hypergraph representations that capture higher-order
relations among knowledge components \cite{thmn2025,hgkt2025,htkt2025} and
dynamic, temporally-evolving graph structure \cite{dygkt2024}.

\todo{Full introduction to be written once the Tier 1 quantitative results
exist (docs/execution-plan.md M5). Draft the gap/motivation/contribution
narrative from docs/idea-D-plan.md sections 1-2 below as a starting
skeleton; do not submit with placeholder contribution claims.}

\subsection{Motivation}
Three concrete gaps motivate this work:

\begin{itemize}
  \item \textbf{Heterogeneous relations are absent from graph-KT.}
  Hypergraph-KT models \cite{thmn2025,hgkt2025,htkt2025} group multiple
  knowledge components into a single hyperedge, but the hyperedge always
  contains only KC nodes. Teacher intervention, hint content, peer
  discussion, and instructional video -- entities with documented influence
  on learning outcomes -- are not represented as graph nodes in any
  published graph-KT model we are aware of.
  \item \textbf{Graph provenance is rarely audited.} \cite{daominh2026p0}
  shows that pooling a knowledge-component graph from the full interaction
  log before a learner-based train/test split lets held-out information
  reach a graph-augmented KT model even when the sequence optimizer itself
  trains only on the training fold. \cite{daominh2026p0}'s protocol audits
  this channel for \emph{pairwise} prerequisite/similarity edges; no
  published audit protocol currently covers \emph{hyperedges} spanning
  heterogeneous entity types.
  \item \textbf{Causal claims about pedagogical intervention are rare.} Most
  graph-KT models report predictive accuracy only, without attempting to
  separate a hint or teacher action's causal effect on subsequent
  correctness from confounded correlation.
\end{itemize}

\subsection{Contributions}
\begin{itemize}
  \item \textbf{C1.} A heterogeneous hypergraph construction procedure
  (Section~\ref{sec:tier1-graph}) that groups \cite{daominh2026p0}'s
  already leakage-audited, acyclic prerequisite edges into
  concept-prerequisite hyperedges, and defines (interface-complete;
  data-blocked, see Section~\ref{sec:limitations}) session,
  discussion-thread, and teacher-intervention hyperedge types.
  \item \textbf{C2.} A hyperedge-level extension of \cite{daominh2026p0}'s
  four-diagnostic leakage taxonomy, adding three classes --
  group-membership, cross-modal, and intervention-timing leakage -- that a
  pairwise audit structurally cannot detect (Section~\ref{sec:leakage-ext}).
  \item \textbf{C3.} A causal-effect layer (Section~\ref{sec:causal})
  estimating the treatment effect of a hint or teacher intervention via
  inverse-propensity weighting, as a design-target starting point before a
  more elaborate causal architecture is attempted.
  \item \textbf{C4.} Direct reuse of \cite{daominh2026p0}'s published,
  leakage-controlled baseline results (BKT, DKT, AKT, simpleKT, GKT, GIKT,
  SKT, DyGKT, DGEKT on XES3G5M/ASSISTments~2012/Junyi Academy) as the Tier~1
  comparison column, avoiding redundant baseline retraining under matched
  protocol and split seeds.
  \item \textbf{C5.} A deliberately-scoped Tier~2 pilot (AgentDG-KT,
  Section~\ref{sec:tier2}) demonstrating that an LLM-agent interpretability
  layer reading Tier~1's frozen output -- rather than re-inferring knowledge
  state from scratch -- is tractable on a single consumer GPU, gated by a
  mandatory Critic agent.
\end{itemize}

\todo{Once results exist, add a C6 sentence stating the headline empirical
finding (or its absence/negative result) -- P0's own Introduction models
this pattern (see its C5) and it should be mirrored here honestly, including
if Tier 1's hypergraph gain turns out to be small or backbone-specific.}

% =============================================================================
\section{Related Work}
% =============================================================================

\subsection{Graph- and hypergraph-based knowledge tracing}
\todo{Expand from docs/idea-D-plan.md and Phien ban A section 1 comparison
table. Cite nakagawa2019gkt, yang2021gikt, dygkt2024, thmn2025, hgkt2025,
htkt2025, psykt2024, hhskt2023.}

\subsection{LLM agents in education and graph-augmented LLM agents}
\todo{Cite agentcat2026, llmedu2025, graphaugmentedllm2025, graphrag2025.
Distinguish AgentDG-KT's ``agent reads frozen Tier-1 output'' design from
prior work where an LLM agent infers knowledge state end-to-end.}

\subsection{Leakage-controlled evaluation in knowledge tracing}
\todo{This section should summarize daominh2026p0 in enough detail that a
reader unfamiliar with the companion paper can still follow
Sections~\ref{sec:tier1-graph}--\ref{sec:leakage-ext} of this paper.}

\subsection{Causal inference in educational data mining}
\todo{Position the propensity-score layer (Section~\ref{sec:causal})
against prior causal-KT / causal-EDM work once identified during the
Related Work pass.}

% =============================================================================
\section{Preliminaries}
% =============================================================================

Let $\mathcal{D} = \{(u, t, q, c, y)\}$ be a KT interaction log with learner
$u$, timestamp $t$, item $q$, knowledge component $c$, and correctness label
$y \in \{0, 1\}$, following \cite{daominh2026p0}'s notation. A learner-based
temporal split $\mathcal{F} = (\mathcal{F}_{\mathrm{train}},
\mathcal{F}_{\mathrm{val}}, \mathcal{F}_{\mathrm{test}})$ assigns each
learner to exactly one fold (ratios $0.7/0.1/0.2$, three folds with split
seeds $42$/$43$/$44$, matching \cite{daominh2026p0}'s protocol so results
remain directly comparable, see Section~\ref{sec:reuse}).

We extend this notation with a \emph{hyperedge} $h = (\kappa,
\{(\tau_i, e_i)\}_{i=1}^{k}, f)$: a kind label $\kappa \in
\{\text{concept-prerequisite}, \text{session}, \text{discussion-thread},
\text{teacher-intervention}\}$, a set of $k$ typed member entities
$(\tau_i, e_i)$ with entity type $\tau_i \in \{\text{student, exercise,
concept, teacher, hint, forum, discussion, video}\}$, and a fold identifier
$f$. A heterogeneous hypergraph $\mathcal{H}_f = (\mathcal{V}_f,
\mathcal{E}_f^{\mathrm{hyper}})$ for fold $f$ is a set of such hyperedges
over the typed node set $\mathcal{V}_f$.

\subsection{Design goal}
For each fold $f$, the goal is a hypergraph $\mathcal{H}_f$ constructed
using only $\mathcal{F}_f^{\mathrm{train}}$ evidence, audited for the
leakage classes in Table~\ref{tab:leakage-taxonomy}
(Section~\ref{sec:leakage-ext}), that a downstream Tier~1 encoder consumes
to predict $y$ for held-out interactions.

% =============================================================================
\section{Tier 1: DH$^2$-KT}
% =============================================================================

\subsection{Heterogeneous hypergraph construction}
\label{sec:tier1-graph}

\todo{Describe build\_concept\_prerequisite\_hyperedges() precisely
(Algorithm~\ref*{alg:concept-hyperedge} below is a first draft matching
dh2a\_kt/hyperedge/construction.py -- keep this in sync with the actual code
as it is implemented per docs/execution-plan.md M1-M2).}

\begin{algorithmic}[1]
\STATE \textbf{Input:} audited, acyclic prerequisite edges
$E_{\mathrm{pre}}^f$ from \cite{daominh2026p0}'s Algorithm~2; chain length
bounds $[\ell_{\min}, \ell_{\max}]$
\STATE \textbf{Output:} concept-prerequisite hyperedge set
$\mathcal{E}^{\mathrm{cprereq}}_f$
\STATE Build directed graph $G$ from $E_{\mathrm{pre}}^f$
\FOR{each root $r$ with in-degree 0 in $G$}
  \FOR{each simple path $p$ from $r$ with $\ell_{\min} \le |p| \le
  \ell_{\max}$}
    \STATE emit hyperedge with members $\{(\text{concept}, c) : c \in p\}$
  \ENDFOR
\ENDFOR
\end{algorithmic}
\label{alg:concept-hyperedge}

Session, discussion-thread, and teacher-intervention hyperedges are
specified with the same interface but are \emph{data-blocked}: no public KT
benchmark studied by \cite{daominh2026p0} (XES3G5M, ASSISTments~2012, Junyi
Academy) ships Hint, Forum, or Teacher-action logs synchronized to the
interaction stream. \todo{State the final resolution once
docs/execution-plan.md M2 is attempted -- either wired against ES-KT-24
\cite{esk2024} for the Video/Hint case, or explicitly scoped to Tier~2
case-study only (Section~\ref{sec:tier2}), per docs/idea-D-plan.md
section~5.}

\subsection{Hyperedge-level leakage audit}
\label{sec:leakage-ext}

\cite{daominh2026p0} audits five leakage classes at the level of pairwise
edge provenance (structural, edge, target, temporal, cold-start
neighbourhood) using four scalar diagnostics: $\mathrm{ECR}^{\mathrm{flag}}$,
$\mathrm{ECR}^{\mathrm{overlap}}$, $|\rho|$, and $\mathrm{TBMR}$. These
diagnostics are reused unmodified (Section~\ref{sec:reuse}) on the pairwise
projection of a hyperedge set. Table~\ref{tab:leakage-taxonomy} adds three
classes that a pairwise audit cannot express because they require a
\emph{group} of heterogeneous members, not a single edge.

\begin{table}[!t]
\caption{Hyperedge-level leakage taxonomy (extends \cite{daominh2026p0}
Table~1)}
\label{tab:leakage-taxonomy}
\centering
\begin{tabular}{@{}p{2.1cm}p{2.9cm}p{2.4cm}@{}}
\toprule
Class & Source & Diagnostic \\
\midrule
Group-membership & A hyperedge with $\ge 1$ member drawn from a held-out
interaction while other members are train-only & $\mathrm{compute\_group\_}$ $\mathrm{membership\_}$ $\mathrm{leakage}$ \\
Cross-modal & A Forum/Video embedding pooled over content timestamped after
the prediction cutoff & $\mathrm{compute\_cross\_}$ $\mathrm{modal\_leakage}$ \\
Intervention-timing & A Teacher/Hint hyperedge whose content was effectively
chosen using knowledge of the student's later correctness & $\mathrm{compute\_intervention\_}$ $\mathrm{timing\_leakage}$ \\
\bottomrule
\end{tabular}
\end{table}

\todo{Report actual leakage rates here once docs/execution-plan.md M1 runs
on real XES3G5M fold data -- Table~\ref{tab:leakage-taxonomy} currently
documents the diagnostic definitions only, not measured values.}

\subsection{Encoding and hypergraph attention}
\todo{Section pending docs/execution-plan.md M3 (dh2a\_kt/models/dh2\_kt.py
DH2KT.forward() implementation). Do not describe the attention mechanism in
past tense until it exists and has passed the sanity overfit test in M3's
Definition of Done.}

\subsection{Causal-effect layer}
\label{sec:causal}
We estimate the average treatment effect (ATE) of a binary intervention
(e.g., hint given) on next-step correctness via inverse-propensity
weighting, following the identification assumption that intervention
assignment is unconfounded given observed features -- a strong assumption on
observational KT data that we report as a limitation rather than a proven
property (Section~\ref{sec:limitations}). Given confounders $X$, treatment
$T \in \{0,1\}$, and outcome $Y$, the propensity score $\hat{e}(X) =
P(T=1\mid X)$ is estimated by logistic regression, and
\begin{equation}
\widehat{\mathrm{ATE}} = \frac{\sum_i T_i Y_i / \hat{e}(X_i)}{\sum_i T_i /
\hat{e}(X_i)} - \frac{\sum_i (1-T_i) Y_i / (1-\hat{e}(X_i))}{\sum_i (1-T_i) /
(1-\hat{e}(X_i))}.
\label{eq:ate}
\end{equation}
\todo{Report propensity-overlap diagnostics alongside any ATE number, per
the warning already implemented in dh2a\_kt/models/causal\_layer.py -- an
ATE without an overlap check is not publishable on its own.}

% =============================================================================
\section{Tier 2: AgentDG-KT pilot}
\label{sec:tier2}
% =============================================================================

Tier~2 is deliberately scoped as a pilot, not a benchmark at Tier~1's
quantitative rigor (docs/idea-D-plan.md, following the feasibility analysis
in the accompanying project discussion: a full multi-agent LLM re-inference
pipeline over millions of interactions is not tractable on a single
consumer GPU within a realistic project timeline). Three agents operate on
Tier~1's \emph{frozen} output rather than re-inferring knowledge state:

\begin{itemize}
  \item \textbf{Diagnostician Agent} produces a natural-language rationale
  for a given Tier-1 prediction, explicitly instructed not to override the
  given probability.
  \item \textbf{Tutor/Hint Agent} proposes the next hint/exercise and writes
  a new session hyperedge back into the graph (the agentic closed loop).
  \item \textbf{Critic Agent} is a \emph{mandatory} quality gate that flags
  any Diagnostician output contradicting or ignoring the Tier-1 number it
  was meant to explain, before anything is treated as ground truth in a
  case study.
\end{itemize}

\todo{Report pilot results (sample size, hallucination/flag rate, human
evaluation if collected) here once docs/execution-plan.md M8-M9 complete.
Do not report Tier 2 metrics using Tier 1's AUC/ACC framing -- the two tiers
use different evaluation standards by design (Section~\ref{sec:eval}).}

% =============================================================================
\section{Experimental setup}
% =============================================================================

\subsection{Datasets}
\label{sec:reuse}
Following \cite{daominh2026p0}'s benchmark-role assignment, XES3G5M
\cite{liu2023xes3g5m} (18{,}066 learners, 7{,}652 items, 865 KCs,
6{,}413{,}353 interactions) is the \emph{primary} benchmark: low
consecutive-KC-repeat rate, non-saturated headline AUC, and the clearest
graph-backbone separation among the corpora \cite{daominh2026p0} studied.
ASSISTments~2012 is secondary (single-skill Q-matrix; leakage-ablation
sanity). Junyi Academy provides the ground-truth cross-validation role via
its expert-curated prerequisite annotation.

\subsection{Baselines reused from P0}
We reuse \cite{daominh2026p0}'s published baseline results
(\texttt{baselines/p0\_reused/}, pinned to commit
\texttt{5e3d6a0}\todo{confirm final vendored commit at submission time}) for
BKT, DKT, AKT, simpleKT, GKT, GIKT, SKT, DyGKT, and DGEKT under train-only,
leakage-controlled graph construction, rather than retraining these nine
models. Reuse is valid only if our preprocessing exactly matches P0's (same
hashed \texttt{kc\_id} space, same Q-matrix); \todo{report the verification
performed via scripts/01\_verify\_p0\_preprocessing\_match.py here.}

\subsection{Manipulation check}
Following \cite{daominh2026p0}'s Section~4.7 reasoning, we do not interpret
any ablation result as evidence of hypergraph benefit unless DH$^2$-KT first
passes a manipulation check: AUC must move meaningfully under near-total
($p=0.90$) destruction of the constructed hyperedge set. A backbone whose
AUC does not move under total destruction is \emph{graph-inert}, and its
ablation numbers carry no interpretable meaning about graph benefit.
\todo{Report pass/fail and $\Delta$AUC here once
docs/execution-plan.md M6 runs.}

\label{sec:eval}

% =============================================================================
\section{Results}
% =============================================================================
\todo{This entire section is empty by design (docs/execution-plan.md
M5-M9 not yet run at the time this skeleton was drafted). Populate with:
(1) Tier-1 AUC/ACC/NLL vs.\ reused P0 baselines, with 95\% CI; (2) ablation
table (per-node-type, hypergraph-vs-pairwise, static-vs-dynamic,
with/without causal layer); (3) manipulation-check result; (4) ground-truth
cross-validation F1 on Junyi; (5) Tier-2 pilot qualitative table. Do not
fabricate placeholder numbers -- leave this section unwritten until real
results exist, per the anti-pattern warned against in every prior planning
document for this project.}

% =============================================================================
\section{Discussion}
% =============================================================================

\subsection{Threats to validity}
\todo{Adapt from docs/idea-D-plan.md section 10 risk table once results
exist to discuss.}

\subsection{Limitations}
\label{sec:limitations}
\begin{itemize}
  \item Session, discussion-thread, and teacher-intervention hyperedges
  remain data-blocked on all three core benchmarks; Tier~1's quantitative
  claims are scoped to Student-Exercise-Concept(-Hint-Video) unless
  Section~\ref{sec:tier1-graph}'s ES-KT-24 integration is completed.
  \item The causal-effect layer's identification assumption (no unmeasured
  confounding given observed features) is stated, not proven, on
  observational KT data.
  \item The hyperedge leakage audit's pairwise-projection approximation
  (Section~\ref{sec:leakage-ext}) weights all projected pairs equally; a
  hyperedge-native statistic is future work.
\end{itemize}

% =============================================================================
\section{Conclusion and future work}
% =============================================================================
\todo{Write after Results/Discussion are complete.}

\section*{Acknowledgment}
\todo{Funding, compute resources (RTX 3090 24GB), and any human evaluators
for Tier 2 (Section~\ref{sec:tier2}) to acknowledge.}

\bibliographystyle{IEEEtran}
\bibliography{references}

\end{document}
